\documentclass[9pt,notes=hide]{beamer}

\usepackage{../../latex/custom}
\usepackage{../../latex/jsr}
\usepackage{../../latex/nsi}

\usetheme[progressbar=frametitle,block=fill]{metropolis}
\usepackage{appendixnumberbeamer}
\setbeamertemplate{frametitle continuation}{\hfill(\insertcontinuationcount)} %

\newcommand{\setsection}[1]{
\setbeamertemplate{footline}{
  \begin{beamercolorbox}[wd=\textwidth, sep=3ex]{footline}%
    #1
    \usebeamerfont{page number in head/foot}%
    \usebeamertemplate*{frame footer}
    \hfill%
    \usebeamertemplate*{frame numbering}
  \end{beamercolorbox}%
}
}

\usepackage{csquotes}
\usepackage{tikz}
\usetikzlibrary{automata, positioning, arrows, patterns}
\usetikzlibrary{matrix,decorations.pathreplacing,chains,fit,shapes,calc} %
\usetikzlibrary{calc,decorations.pathmorphing,decorations.markings}

\usepackage{pgfplots}
    \pgfplotsset{compat=1.11}
\usepackage{pgfplotsthemetol} %
\usepackage{subcaption} %
\usepackage{array}

\newcommand{\abs}[1]{\left| #1 \right|} %
\newcommand{\Init}{\mathrm{Init}} %
\newcommand{\comment}[1]{{\color{red}\textbf{{#1}}}}

\usepackage[super]{nth}

\newcommand{\expe}[2]{\la #1, #2 \ra}
\newcommand{\piB}{\pi_{\Image(B)^\perp}}
\DeclareMathOperator{\gauge}{g}

\newcommand\Wider[2][3em]{%
\makebox[\linewidth][c]{%
  \begin{minipage}{\dimexpr\textwidth+#1\relax}
  \raggedright#2
  \end{minipage}%
  }%
}


\title{Set Programming: Theory and Computation}
\subtitle{Private Defense}
\date{April 28, 2020}
\author{Beno\^{i}t \textsc{Legat}}
\institute{Jury: Rapha\"{e}l \textsc{Jungers}, Pablo \textsc{Parrilo}, Yurii \textsc{Nesterov}, Anthony \textsc{Papavasiliou}, Paulo \textsc{Tabuada}, Roland \textsc{Keunings}}

\begin{document}

\maketitle




\begin{frame}{Table of contents}
  \tableofcontents
\end{frame}

\setsection{\S4}

\begin{frame}{Set Program}
\begin{equation*}
  \begin{array}{rl}
    \min\limits_{\mathcal{S}_1 \subseteq \R^{n_1}, \ldots, \mathcal{S}_N \subseteq \R^{n_N}} & f(\mathcal{S}_1, \ldots, \mathcal{S}_N)\\
    \text{subject to:} & g_i(\mathcal{S}_1, \ldots, \mathcal{S}_N) \subseteq h_i(\mathcal{S}_1, \ldots, \mathcal{S}_N), \quad i=1,2,\ldots,M.
  \end{array}
\end{equation*}

\begin{block}{Linear inclusion}
\begin{align*}
  A\mathcal{S}_i \subseteq \mathcal{S}_j,\\
  \mathcal{S}_i \subseteq A^{-1}\mathcal{S}_j.
\end{align*}
\end{block}
\begin{block}{Algebraic inclusion}
\begin{align*}
  \mathcal{S}_i & \subseteq \pi_{[n_i], n_i + m}
  \begin{bmatrix}
    A & B
  \end{bmatrix}^{-1} \mathcal{S}_j,\\
  C\mathcal{S}_i & \subseteq E\mathcal{S}_j.
\end{align*}
\end{block}
\end{frame}

\section{Guarantees}

\setsection{\S4.4}

\begin{frame}{L\"{o}wner-John ellipsoid of intersection of ellipsoids}
  John ellipsoid:
  \begin{equation*}
    \begin{array}{rl}
      \max_{\mathcal{E}_P \subseteq \R^n} & \vol(\mathcal{E}_P)\\
      \text{subject to:} & \mathcal{E}_P \subseteq \mathcal{C}.
    \end{array}
  \end{equation*}
  \begin{columns}
    \begin{column}{0.5\textwidth}
      \textbf{Theorem}:
      \alert{Unique} $\opt{\mathcal{E}_P}$ and ($\mathcal{C}$ symmetric):
      \[ \mathcal{E}_P^\star \subseteq \mathcal{C} \subseteq \sqrt{n} \mathcal{E}_P^\star. \]

      \textbf{Conjecture} (\alert{Proved} for aligned axes or $m = 2$):
      $\mathcal{C}$ intersection of $m$ ellipsoids (joint eccentricity $\kappa$, $l = \min(n, m)$):
      \[
        \mathcal{E}_P^\star \subseteq \mathcal{C}
        \subseteq
        \sqrt{\frac{l\kappa^2}{l-1+\kappa^2}} \mathcal{E}_P^\star.
      \]
      If $\kappa \to \infty$, then
      \[
        \mathcal{E}_P^\star \subseteq \mathcal{C}
        \subseteq
        \sqrt{\min(n, m)} \mathcal{E}_P^\star.
      \]
    \end{column}
    \begin{column}{0.5\textwidth}
      \centering
      \begin{tikzpicture}[scale=1]
          \draw (0, 0) circle (1);
          \draw (-1, -1) rectangle (1, 1);
          \draw[dashed] (0, 0) circle (1.4142135623730951);
          \draw[fill] (0, 0) circle (.02);
      \end{tikzpicture}
      \centering
      \begin{tikzpicture}[scale=0.8]
          \draw (0, 0) ellipse (1 and 2);
          \draw (0, 0) ellipse (2 and 1);
          \draw[<->] (0, 0) to (2, 0);
          \node at (1.6, 0.15) {$\kappa$};
          \draw (0, 0) circle (1);
          \draw[dashed] (0, 0) circle (1.2649110640673518);
          \draw[fill] (0, 0) circle (.02);
      \end{tikzpicture}
    \end{column}
  \end{columns}
\end{frame}

\setsection{\S4.2, \S5.1, \S5.2}

\begin{frame}{Linear inclusions and switched systems}
  \begin{columns}
    \begin{column}{0.5\textwidth}
      Matrices $A_1, \ldots, A_m \in \R^{n \times n}$.
      \begin{block}{Switched system}
        \begin{equation*}
          x_k = A_{\sigma_k} x_{k-1}, \quad \sigma_k \in [m]
        \end{equation*}
      \end{block}

      \vspace{-1em}

      \begin{equation*}
        \begin{array}{rl}
          \jsr = & \min_{\mathcal{S} \subseteq \R^{n}}\gamma\\
          \text{subject to:} & \frac{A_\sigma}{\gamma} \mathcal{S} \subseteq \mathcal{S}, \quad i \in [m].
        \end{array}
      \end{equation*}

      \vspace{1em}

      $\ujsrsos$: for polyset $\mathcal{S}$ of degree $2d$.

      \vspace{0.4em}

      \begin{block}{Guarantee}
        \vspace{-1em}
        \begin{align*}
          \jsr & \le \ujsrsos \le \alpha_{2d} \jsr\\
          \alpha_2 & = \sqrt{\min(n, m)}\\
          \alpha_{2d} & = \min\left({n+d-1 \choose d}, m\right)^{\frac{1}{2d}}.
        \end{align*}
      \end{block}
    \end{column}
    \begin{column}{0.5\textwidth}

      Labelled graph $G(V, E)$.
      \begin{block}{Constrained switched system}
        \begin{equation*}
          x_k = A_{\sigma_k} x_{k-1},
          (u_{k-1}, u_k, \sigma_k) \in E
        \end{equation*}
      \end{block}

      \vspace{-2em}

      \begin{equation*}
        \begin{array}{rl}
          \cjsr = & \min_{\mathcal{S}_u \subseteq \R^{n}, u \in V} \gamma\\
          \text{subject to:} & \frac{A_\sigma}{\gamma} \mathcal{S}_u \subseteq \mathcal{S}_v, (u, v, \sigma) \in E.
        \end{array}
      \end{equation*}

      $\jsrsos$: for polysets $\mathcal{S}_u$ of degree $2d$.

      \begin{block}{Guarantee}
        \vspace{-1em}
        \[ \cjsr \le \jsrsos \le \alpha_{2d} \cjsr \]
        \vspace{-2em}
        \begin{align*}
          \alpha_2 & = \sqrt{\min(n, 2^{h(E)})}\\
          \alpha_{2d} & = \min\left({n+d-1 \choose d}, 2^{h(E)}\right)^{\frac{1}{2d}}.
        \end{align*}
      \end{block}
    \end{column}
  \end{columns}
\end{frame}

\setsection{\S5.3}

\begin{frame}{Conic duality}
  Infeasibility certificate $\mu_\sigma$ for polysets $\to$ invalid for generic set program.
  Rounding procedure generates \alert{backward} switching sequence:

  If \alert{atomic}, take $x_0$:
  \[ \sigma_{k+1} = \argmax_\sigma \mu_\sigma(\{x_k/\norm{x_k}_2\}), \quad x_{k+1} = A_{\sigma_{k+1}} x_k \]

  \begin{block}{Guarantee}
    \[ \lim_{k \to \infty} \norm{A_{\sigma_1} \cdots A_{\sigma_k}}^{\frac{1}{k}} = \cjsr = \jsrsos. \]
  \end{block}

  Else, take $p_0(x)$ SOS:
  \[ \sigma_{k+1} = \argmax_\sigma \expe{\mu_\sigma}{p_k(x)}, \quad p_{k+1}(x) = p_k(A_{\sigma_{k+1}} x) \]

  \begin{block}{Guarantee}
    \[ 2^{-\frac{h(E)}{2d}}\jsrsos \leq \lim_{k \to \infty} \norm{A_{\sigma_1} \cdots A_{\sigma_k}}^{\frac{1}{k}} \leq \cjsr \leq \jsrsos. \]
  \end{block}
\end{frame}

\section{Geometric duality}

\setsection{\S4.3}

\begin{frame}{Geometric duality}
  \begin{block}{Discrete-time}
  \begin{align*}
    C\mathcal{S}_i %
    \subseteq E\mathcal{S}_j.
  \end{align*}

  \vspace{-1em}

  Inequality with support function or polar:

  \vspace{-1em}

  \begin{align*}
    \forall z \in \R^r, \quad \supfun{C^\Tr z}{\mathcal{S}_i} & \le \supfun{E^\Tr z}{\mathcal{S}_j}\\
    \forall z \in \R^r, \quad \gauge(\polar{\mathcal{S}_i}, C^\Tr z) & \le \gauge(\polar{\mathcal{S}_j}, E^\Tr z).
  \end{align*}
  \end{block}

  \begin{block}{Continuous time}
  \[ \forall x \in \partial \mathcal{S}_i, C x \in E T_{\mathcal{S}_i}(x). \]

  Inequality with support function or polar:
  \begin{align*}
    \forall z \in \mathbb{R}^r, \la z, C \grad \supfun{E^\Tr z}{\mathcal{S}_i} \ra & \le 0\\
    \forall z \in \mathbb{R}^r, \la z, C \grad \gauge(\polar{\mathcal{S}_i}, E^\Tr z) \ra & \le 0.
  \end{align*}
  \end{block}
\end{frame}

\setsection{Example~40, Example~42}

\begin{frame}{Examples}
  \begin{columns}
    \begin{column}{0.5\textwidth}
      \centering

      \alert{Primal} space

      \includegraphics[width=0.8\textwidth]{SimpleControlledInvariantEllipsoid.png}
      \includegraphics[width=0.8\textwidth]{SimpleControlledInvariantPolySet.png}
    \end{column}
    \begin{column}{0.5\textwidth}
      \centering

      \alert{Dual} space

      \includegraphics[width=0.8\textwidth]{SimpleControlledInvariantPolarEllipsoid.png}
      \includegraphics[width=0.8\textwidth]{SimpleControlledInvariantPolarPolySet.png}
    \end{column}
  \end{columns}
\end{frame}

\setsection{\S6}

\begin{frame}{Controlled invariant set}
  \begin{block}{Discrete-time}
    \[ x_{k+1} = Ax_k + Bu_k. \]

    Controlled invariance of $\mathcal{S}$:
    \[ \forall x \in \mathcal{S}, \exists u, A x + B u \in \mathcal{S}. \]

    \vspace{-1em}

    Equivalent to
    \[ \piB{} \mathcal{S} \subseteq \piB{} A \mathcal{S}. \]
  \end{block}

  \begin{block}{Continuous time}
    \[ \dot{x}(t) = Ax(t) + Bu(t). \]

    Controlled invariance of $\mathcal{S}$:
    \[ \forall x \in \partial \mathcal{S}, \exists u, A x + B u \in T_\mathcal{S}(x). \]

    \vspace{-1em}

    Equivalent to
    \[ \forall x \in \partial \mathcal{S}, \piB{} A x \in \piB{} T_\mathcal{S}(x). \]
  \end{block}
\end{frame}

\section{Applications}

\setsection{\S6.3}

\begin{frame}{Model Predictive Control}
  \begin{center}
    \begin{tikzpicture}
      \tikzstyle{spring}=[thick,decorate,decoration={zigzag,pre length=0.1cm,post length=0.1cm,segment length=6}]
\tikzstyle{damper}=[thick,decoration={markings,
  mark connection node=dmp,
  mark=at position 0.5 with
  {
    \node (dmp) [thick,inner sep=0pt,transform shape,rotate=-90,minimum width=10pt,minimum height=2pt,draw=none] {};
    \draw [thick] ($(dmp.north east)+(2pt,0)$) -- (dmp.south east) -- (dmp.south west) -- ($(dmp.north west)+(2pt,0)$);
    \draw [thick] ($(dmp.north)+(0,-4pt)$) -- ($(dmp.north)+(0,4pt)$);
  }
}, decorate]
\draw[thick] (0, 0) -- (2, 0) -- (2, 1) -- (0.5, 1) -- (0, 0.5) -- (0, 0);
\draw [spring] (2, 0.8) -- (3, 0.8);
\draw [damper] (2, 0.2) -- (3, 0.2);
\draw[thick] (3, 0) rectangle (5, 1);
\draw [spring] (5, 0.8) -- (6, 0.8);
\draw [damper] (5, 0.2) -- (6, 0.2);
\draw[thick] (6, 0) rectangle (8, 1);
\draw[thick] (0.5, 0.1) circle (0.3);
\fill[color=lightgray] (0.5, 0.1) circle (0.3);
\draw[thick] (1.5, 0.1) circle (0.3);
\fill[color=lightgray] (1.5, 0.1) circle (0.3);
\draw[thick] (3.5, 0.1) circle (0.3);
\fill[color=lightgray] (3.5, 0.1) circle (0.3);
\draw[thick] (4.5, 0.1) circle (0.3);
\fill[color=lightgray] (4.5, 0.1) circle (0.3);
\draw[thick] (6.5, 0.1) circle (0.3);
\fill[color=lightgray] (6.5, 0.1) circle (0.3);
\draw[thick] (7.5, 0.1) circle (0.3);
\fill[color=lightgray] (7.5, 0.1) circle (0.3);
\node at (1, 0.5) {$v_0$};
\node at (4, 0.5) {$v_1$};
\node at (7, 0.5) {$v_2$};
\node at (2.5, 1.1) {$d_1$};
\node at (5.5, 1.1) {$d_2$};

      \draw[thick] (-1, 0) to (-1, 1);
      \draw[thick,color=red] (-1, 1.3) circle (0.3);
      \node at (-1, 1.3) {\scriptsize \bf 15.6};
    \end{tikzpicture}

    \includegraphics[width=0.45\textwidth]{CruiseControlSets_d_u.png}
    ~
    \includegraphics[width=0.45\textwidth]{CruiseControlSets_v0_u.png}
  \end{center}
\end{frame}

\setsection{\S6.4}

\begin{frame}{Stochastic Programming}
  \begin{columns}
    \begin{column}{0.6\textwidth}
  No relatively complete recourse:
  \begin{itemize}
    \item feasibility cuts or
    \item penalization.
  \end{itemize}
  $\to$ Slow down Stochastic Dual Dynamic Programming.

  Feasibility part of stochastic program $\to$ hybrid control algebraic system.

  Controlled invariant set ensures relatively complete recourse.
    \end{column}
    \begin{column}{0.4\textwidth}
    \includegraphics[width=\textwidth]{reservoir.jpg}~%
    \end{column}
  \end{columns}

\end{frame}

\setsection{\S7}

\begin{frame}{Information Theory}
  Invariant under inner- and self-adhesivity.
  \begin{align*}
    \ia{\Jis}{\Kis}{\Iis}(h) & = \{\, g \in \Es[n'] \mid \forall \Iis \subseteq \Lis \subseteq \Jis \cap \Kis, g_{\Lis} = h_{\Lis \cap \Jis} + h_{\Lis \cap \Kis} - h_{\Iis} \,\}, & \Iis & = \Jis \cap \Kis,\\
    \sa{\Jis}{\Iis}(h) & = \{\, g \in \Es[n'] \mid \forall \Iis \subseteq \Lis \subseteq \Jis' \cap \Kis, g_{\Lis} = h_{\Lis \cap \Jis'} + h_{\Lis \cap \Kis} - h_{\Iis} \,\}, & \Iis & \subseteq \Jis,
  \end{align*}

  Generate feasibility cuts with Dual Dynamic Programming.

  \includegraphics[width=0.33\textwidth]{tetra0pm.png}~%
  \includegraphics[width=0.33\textwidth]{tetra0zy.png}~%
  \includegraphics[trim=6cm 2cm 6cm 7cm,clip,width=0.33\textheight]{tetra0bound15789.png}

\end{frame}


\setsection{}

\begin{frame}{Conclusion}
  \begin{itemize}
    \item Set programs allow to model many important problems.
    \item \alert{Guarantees} for templates can be developed for subclasses.
    \item Infeasibility certificate for a template may be \alert{rounded}.
    \item Geometric \alert{similarities} of different templates.
    \item \alert{Geometric duality} enable reformulation in functional inequality.
  \end{itemize}
  \begin{center}
    \includegraphics[width=0.5\textwidth]{FinConjCounterEx.png}
  \end{center}
\end{frame}

\end{document}
